Many machine learning models are based on supervised learning, where the data is
labelled. In the field of computer vision, the labelling of image data is
foundational to the success of machine learning models, as it allows the models
to learn the relationship between the input images and the output labels. To
train such models, a large amount of labelled data is required. In many cases,
the image has a ground truth label which is unknown. In such cases labels can be
provided by human annotators, for example in the context of medical imaging,
wildlife monitoring, and quality control in manufacturing. However, human
labelling is not infallible and is subject to various uncertainties that can
affect the performance of machine learning algorithms. Often, the label assigned
to an image is determined by a majority vote among multiple annotators. This
loses a lot of information about the labelling process and the uncertainty
associated with the labels. Using the methods proposed in this thesis, we can
quantify the uncertainty in the labelling process and incorporate this
information in future work into the training of machine learning models to
improve their performance, as well as to provide a measure of confidence in the
predictions made by the models.\bigskip

This thesis addresses the issue of quantifying labelling uncertainty in image
data, specifically focusing on the widely-used \textit{CIFAR-10} dataset,
which contains images with their respective true labels. We will use an extended
version of the \textit{CIFAR-10} dataset, called the \textit{CIFAR-10H} dataset,
that was introduced in \cite{petersonHumanUncertaintyMakes2019} and extends the
\textit{CIFAR-10} dataset by providing human labels. The goal is to quantify the
uncertainty in the labelling process and to compare the results obtained under
the assumption of missing the true labels to the results obtained with the true
labels, and thus verifying the functionality of the proposed method. For the
methodology we build upon, we refer to
\cite{hechingerCategorisingWorldLocal2024} and
\cite{mclachlanBayesianApproachMixture2000} 
who use a finite mixture model to
estimate the uncertainty in the labelling process. We also extend the previous
work by considering the time taken by the annotators to label the images as a
source of labelling uncertainty.\bigskip

The thesis is structured as follows. In Chapter \ref{Data}, we introduce the
\textit{CIFAR-10} and \textit{CIFAR-10H} datasets and describe the labelling
process. In Chapter \ref{uncertainty_estimation}, we present the methodology
used to estimate the labelling uncertainty. This includes the Bayesian latent
class mixture model used to estimate the uncertainty in the labelling process,
as well as the stochastic expectation-maximization algorithm used to fit the
model. In Chapter \ref{uncertainty_sources}, we show the results of the
uncertainty estimation and compare the results obtained under the assumption of
missing the true labels to the results obtained with the true labels. We also
analyse different sources of labelling uncertainty, such as the subcategories of
the images and the time taken by the annotators to label the images. Finally, in
Chapter \ref{conclusion}, we conclude with a summary and discussion of the
results, limitations of the proposed method, as well as an outlook on possible
future research questions.